\documentclass[12pt]{article}
\usepackage[utf8]{inputenc}
\usepackage[T1]{fontenc}
\usepackage{hyperref}
\title{Wardaman Language Resources for K-12 Curriculum Development}
\author{Compiled by Claude Sonnet 5.5 (claude-sonnet-5-5) for Andr\'as Kornai}
\date{}
\begin{document}
\maketitle

\section{Introduction}
Wardaman (ISO 639-3 wrr; Glottocode ward1246) is a severely endangered Yangmanic language of the Northern Territory, Australia, in country southwest of Katherine. The ANU project describes only two elderly first-language speakers, a brother and sister, at the time of its recordings.

\section{Primary Educational Resources}
No school primers, textbooks or programmes were found in search results.

\section{Linguistic and Reference Materials}
\begin{itemize}
\item Francesca Merlan, \emph{A Grammar of Wardaman} (Mouton de Gruyter, 1994). It has phonology, morphology and syntax, thirty interlinear texts, a semantic-domain vocabulary and a wordlist. Catalogue record: \url{https://search.worldcat.org/oclc/28926390}. Publisher page: \url{https://www.degruyterbrill.com/document/doi/10.1515/9783110871371/html}
\item Project page, ``Wardaman Dictionary, Narrative, Song and Country'' (ANU): \url{https://researchportalplus.anu.edu.au/en/projects/wardaman-dictionary-narrative-song-and-country/}. It aimed to expand a wordlist into a dictionary. I found no standalone published dictionary.
\end{itemize}

\section{Religious and Community Resources}
No scripture or community organisation pages were found. Wikipedia has a page on the Wardaman people: \url{https://en.wikipedia.org/wiki/Wardaman_people}

\section{Digital and Audio Resources}
The ELDP archive holds the collection ``Wardaman Dictionary, Narrative, Song and Country'' (collection 0696), deposited by Francesca Merlan: \url{https://eldp.access.preservica.com/dk0696/}. The collection includes recordings. Access conditions were not checked. The Endangered Languages Project profile is at \url{https://endangeredlanguages.com/lang/2580}.

\section{Curriculum Development Context}
The language has a full grammar and an archived documentation collection, but no known teaching materials. A curriculum would have to be developed with the Wardaman community, using the archive and grammar as sources.
\end{document}